\chapter*{Judas}

So, Judas from the town of Kerioth. The only Judean among 12 Galileans. The relationship between these two Palestinian peoples was not a particularly warm one; this is often cited as an explanation for the less-than-enthusiastic reception the Galilean Jesus enjoyed in Judea and its capital, Jerusalem (``Are you also from Galilee? Search and look, for no prophet has arisen out of Galilee''---Jn 7:52). However, this state of affairs did not stop Judas, a relative latecomer to Christ's sect (if not one of his last followers), from being chosen as one of his 12 Apostles and even becoming their treasurer. This alone is a testament to the level of trust and authority that he enjoyed among Christ's community of believers. And for what it's worth, Jesus doesn't give me the impression of being a holy fool, unable to put two and two together when it comes to worldly matters; on the contrary, he seems quite discerning when it comes to people, not unlike Tsar Fyodor Ioannovich\footnote{A 16\textsuperscript{th}-century Russian tsar whose life was dramatized in Aleksey Tolstoy's play \emph{Tsar Fyodor Ioannovich}; the play depicts him as a deeply moral man who is nonetheless incapable of ruling the country, a task instead taken up by his brother-in-law, Boris Godunov. (Z.B.)} (or, for that matter, Bulgakov's Yeshua).

For good reason, many people have been left unconvinced by the canonical narrative of Judas' ``betrayal for thirty pieces of silver,'' and have sought other explanations for his actions; in this sense, he has certainly become the most popular figure from the Gospel. The alternative explanations that have been proposed vary incredibly widely: they include seething resentment at Christ the ``deceiver,'' whose kingdom, as it suddenly turned out, would not be of this world; a wish to see if the man who claimed to be the Messiah could even save himself; and a desire to hasten the establishment of the Kingdom of God on Earth (or perhaps provoke a popular uprising). Case in point: Franco Zeffirelli's spectacular film \emph{Jesus of Nazareth} consists, for the most part, of stunning ``moving illustrations'' of the Gospel text, yet even here, Judas' subplot is out of line with the canonical version.

Particularly notable in this regard is a heretical version expounded by Borges, which he attributes to the fictional Swedish theologian Nils Runeberg:

\begin{smallquote}
    ...He begins by pointing out how superfluous was the act of Judas. [...] In order to identify a master who daily preached in the synagogue and who performed miracles before gatherings of thousands, the treachery of an apostle is not necessary.\footnote{``Three Versions of Judas'' from \emph{Fictions}, trans. Anthony Kerrigan, p. 139.}
\end{smallquote}
The real motivation for Judas' actions was ``an extravagant and even limitless
asceticism'':

\begin{smallquote}
    The ascetic, for the greater glory of God, degrades and mortifies the flesh; Judas did the same with the spirit. He renounced honor, good, peace, the Kingdom of Heaven, as others, less heroically, renounced pleasure.\footnote{``Three Versions of Judas,'' pp. 140-141}
\end{smallquote}
Hmm... Well, perhaps we won't be soaring to such heavenly heights of theological thought. For now, though, let's focus on finding motivations for Judas' actions that run a little closer to ``the objective reality given us in our sensations.''\footnote{Vladimir Lenin, \emph{Materialism and Empirio-criticism}, trans. Abraham Fineberg, p. 171.}

``Thirty pieces of silver'' as a motivation for betrayal, however, fails to hold water for the most pragmatic reason possible: what would such a paltry sum mean to Judas compared to the possibilities at his disposal as the Apostles' treasurer? If the driving force behind Judas' actions was greed, then it would have suited him far better to quietly keep skimming money from the community's money box. Only a complete idiot (or a sovok\footnote{A sovok (otherwise known as \emph{Homo sovieticus}) is a person who slavishly and naively adheres to a Soviet-era mentality; in Russian, the word literally means ``dustpan.'' (Z.B.)}) would kill the goose that lays the golden egg.

But did Judas actually steal any money? John says so with total confidence (Jn 12:6); it is strange, however, that not one of the Synoptic Evangelists mention such a colorful detail, which would really bring the image of the traitor to life. One can only assume that John, as he had already done before, managed to get things completely mixed up. One could suppose that sometime before the tragedy between Jesus and Judas, there was an intense argument about finances. However, there was no proof that Judas was behind the shortfall; otherwise, Peter and the other Apostles would have found this episode far too significant to exclude from their narratives. Let's keep this in the back of our minds.

I should note that the text of the New Testament contains a direct indication that Christ foresaw Judas' betrayal long before his last trip to Jerusalem. I am talking about John's \emph{interpretation} of the statements his teacher made after giving his Bread of Life Discourse and subsequently having many of his companions desert him:

\begin{smallquote}
    Then Jesus said to the twelve, ``Do you also want to go away?'' But Simon Peter answered Him, ``Lord, to whom shall we go?'' [...] Jesus answered them, ``Did I not choose you, the twelve, and one of you is a devil?'' He spoke of Judas Iscariot, the son of Simon, for it was he who would betray Him, being one of the twelve. (Jn 6:67-71)
\end{smallquote}
However, this is not as simple and clearcut as John would have it.

Let's start with the fact that, generally speaking, it does not follow from this text that Jesus specifically had Judas in mind; while this speculation may be logical in retrospect, it is speculation nonetheless. Furthermore, if one assumes that John's interpretation was valid, a serious theological issue arises: Holy Scripture now explicitly affirms the rather pagan idea that man's actions are fatally predetermined, which seems completely incompatible with all Christian philosophy. So, let's suppose that the statement John recorded was a legitimate one, and let's also suppose that it related directly to Judas. Even then, however, there's nothing that implies that Christ is referring to his \emph{future betrayal}, as opposed to some other act of Judas'---perhaps one that he had just committed at that very moment.

While we're at it, what were the circumstances of Judas' death? The most common version is that of Matthew---he ``went and hanged himself'' (Mt 27:5). Meanwhile, Acts of the Apostles says something completely different:

\begin{smallquote}
    He burst open in the middle and all his entrails gushed out. And it became known to all those dwelling in Jerusalem. (Ac 1:18-19)
\end{smallquote}
Farrar opines, however, that this version is ``not irreconcilable with the first,''\footnote{\emph{The Life of Christ}, p. 419} and even comes up with a sort of hybrid: the rope Judas is hanging from snaps, he hits the ground from a great height, and his belly bursts open... I feel no need to comment on the plausibility of this theory.\footnote{Archpriest Alexander Men was at least honest when he wrote that ``the information contained in Mt 27:3-9 and Ac 1:16-20 is still [?!---K.E.] quite difficult to reconcile.'' (K.E.)}

And what's more: the prospect of Judas committing suicide out of regret, when considered in light of his actions during his teacher's arrest, is completely unconvincing from a psychological perspective. Such a thing has been known to happen with people who commit traitorous acts out of coercion or blackmail; Judas, however, acted on his own initiative, deliberately and in cold blood. We could, of course, accept the canonical version, that at the moment of his betrayal, he was ``possessed by a demon'' and subsequently became deluded. For now, though, let's only turn to this ``explanation'' as a very last resort.

Most puzzling of all, however, is the scene of Christ's arrest in the Garden of Gethsemane. That is, it's not that the number of inconsistencies (assuming we go by the canonical version) is merely large; it's that the entire episode consists of literally nothing but inconsistencies, down to the very last detail. Let me once again yield the floor to Dombrovsky:

\begin{smallquote}
    ``As it happens, there's a great confusion behind the whole story. After all, Christ didn't hide himself away, he spoke in public. He could have been seized perfectly easily any day even without the help of Judas. `Why these swords and staves,' he said when he was arrested. `You've seen me every day, and I preached among you. Why didn't you take me then?''' ``Logical,'' smiled Surovtsev [an NKVD officer --- K.E.]. ``I mean, logical only for Christ. People who are arrested often ask that question. They are ignorant of the operative considerations involved.''\footnote{\emph{The Faculty of Useless Knowledge}, p. 264}
\end{smallquote}

And that right there is exactly what we are going to discuss---just what kinds of ``operative considerations'' were there? Let's ask ourselves a couple of questions.
 
1. Christ's arrest could have taken place at any moment of the day on the streets of Jerusalem; he traveled with very few companions, and they certainly wouldn't have been able to put up any resistance against the temple guards. Recall that the authorities encountered practically no problems when arresting the incomparably more popular John the Baptist. Why in the world did they let the Galileans leave the city and slip into the forested outskirts of Jerusalem, where it would be hard to monitor their movements even in broad daylight? In other words: what exactly did the Judean authorities have to gain by drastically complicating the arrest and conducting it in a quiet, isolated location in the dead of night?

2. Judas' behavior during the arrest seems completely irrational. His task, which he completed successfully, was to lead the arrest team to the site (which was unknown to anyone in Jerusalem). At this point, any normal traitor would try to slip into the shadows (both figuratively and literally), and not thrust himself into the limelight to flaunt his exceptional merits. So why, pray tell, did Judas feel the need to point Christ out in such a public, theatrical fashion? During the time he spent preaching in Jerusalem, Judean secret police operatives certainly managed to get the best ``visual'' possible on him---and the arrest certainly couldn't have gone on without them. And on a related note: why did Judas need to pass himself off as the commander of the arrest team (``concerning Judas, who became a guide to those who arrested Jesus''---Ac 1:16) when in reality, he most certainly wasn't? After all, even if the high priests lost their minds and chose a defector to lead the unit (like they say: once a traitor, always a traitor, no matter who he betrayed), never in their lives would the Roman legionnaires willingly take orders from some Jewish thug who'd just betrayed his boss.

3. While we're at it, it's worth bringing our attention to the composition of the arrest team, which is bizarre from any point of view. First, as was previously mentioned, it consisted not only of Jews, but also of Roman soldiers, who the initiators of the arrest---the high priests---had no direct authority over. \footnote{``The Johannine version that Jesus was arrested by the whole or a part of a Roman cohort, commanded by its tribune, with the Jewish temple police and its commanding officer in attendance, is now being accepted as the true statement of facts by most contemporary scholars'' (Haim Cohn, \emph{The Trial and Death of Jesus}, p. 78). Cohn himself, whose book specifically analyzes the legal intricacies of the ``trial of Jesus of Nazareth,'' considers this detail extremely strange and demanding of special explanation. Small wonder---you have to think the former Attorney General of Israel would be well-informed on how to actually carry out an arrest... (K.E.)} Seeking the legionnaires' participation in the operation would have required them, at the very minimum, to inform the procurator accordingly, but that would mean wasting precious time on the inevitable negotiations---and what for? It could be justified if Caiaphas were expecting serious armed resistance; such a risk was negligible, however---compared to the very real possibility that the alarmed sect might simply vanish into the night, effectively becoming a needle in a haystack. Second, the Romans, who numbered no more than two to three dozen soldiers, were commanded by a ``thousandman'' (a military tribune). The participation in this operation of such a high-ranking officer (the equivalent of a colonel in today's military) clearly speaks to how seriously Pilate took the high priests' ``request for international aid.''\footnote{The exact phrase used by the Soviet government to justify its invasion of Afghanistan in 1979. (Z.B.)} So just why does he start playing the fool the next morning, portraying his relationship to the defendant as one of benevolent neutrality? Third, in addition to the temple guards (as well as the inevitable undercover detectives), the Judean contingent was chock full of high priests' servants carrying clubs. If I may ask: what need was there to suddenly call for such a ``total mobilization,'' and what use could this riff-raff have possibly served---other than to step on the actual professionals' toes?

4. It is unclear why the arrest team took absolutely no measures to detain the Apostles. After all, even if the authorities decided to ignore their previous complicity in spreading subversive propaganda, they still put up direct armed resistance the moment Christ was arrested. \footnote{According to the Old Russian version of \emph{The Jewish War} (IX:3), a great number of people were killed during Jesus' arrest. (K.E.)} Nevertheless, there were no consequences when Peter sliced off the ear of the high priests' servant Malchus (Jn 18:10), and the Apostles were allowed to slip away unimpeded. This is especially unclear considering the events that would occur later that night---the three attempts to arrest Peter specifically due to his membership in Christ's inner circle, completely unrelated to his part in the armed confrontation with the guards.

5. As Sherlock Holmes once said, ``the more outre and grotesque an incident is, the more carefully it deserves to be examined, and the very point which appears to complicate a case is, when duly considered and scientifically handled, the one which is most likely to elucidate it.''\footnote{Arthur Conan Doyle, \emph{The Hound of the Baskervilles}, p. 242.} In our case, that detail is... the torches; yes, the very torches the arrest team brought to the site (e.g. Jn 18:3). The issue is that these events took place during the Jewish Passover, which falls on a full moon. The torches might have been useful in total darkness, if the arrest team needed to sweep a cordoned-off section of the garden (though even then, they would do just as much harm as good) or identify detainees. But what purpose did it serve (and who did it benefit) to illuminate a garden that was already flooded with moonlight, leaving the arrest team exposed on their way to a location they were already familiar with?